% ADAPTIVE-FORMULA-FBDD5
\section*{ADAPTIVE-FORMULA-FBDD5: adaptive clause-query DAGs and exact evasiveness of full $D(4)$}

\textbf{STATUS: PROVED/EXACT for the finite query-depth statements and $D(4)$ evasiveness; COMPUTATIONAL upper bounds for FBDD size. P versus NP remains OPEN.}

\paragraph{Model.}
The input bits are clause-presence bits. An FBDD is a deterministic DAG in which each state may choose its next queried clause adaptively, but no clause bit is queried twice on one source--sink path. Identical residual SAT/UNSAT functions may merge. Thus
\[
\mathrm{OBDD}\subseteq\mathrm{FBDD}\subseteq\mathrm{general\ BP}.
\]

\paragraph{Canonical residual state.}
After a partial presence-bit history, let $U$ be the truth assignments not yet falsified by selected clauses and $R$ the unread clause bits. Future UNSAT means that future selected clauses cover all of $U$. The exact residual Boolean function is represented by the antichain of unread-clause sets capable of killing each uncovered assignment, after duplicate and subsumed positive clauses are removed.

\paragraph{Adaptive finite results.}
The best constructions found are
\[
\begin{array}{c|r|r|r|r}
\text{family}&M&\text{fixed-priority nodes}&\text{adaptive nodes}&\text{exact adaptive depth}\\
\hline
2\text{-SAT}_3&12&86&70&12\\
\text{Horn}_3&12&32&29&9\\
\text{XOR}_3&14&146&122&14\\
\text{two-triple chain}&16&77&24&16\\
D(4)\text{ full}&32&2438&2432&32.
\end{array}
\]
The adaptive sizes are upper bounds, not optimum FBDD sizes.

The 24-node chain FBDD was checked on all
\[
2^{16}=65{,}536
\]
clause-presence inputs with zero mismatches. Its structure is especially simple: the eight clauses of the first local triple occur once globally and the eight clauses of the second triple occur twice globally. Thus adaptive sharing recognizes bounded-width locality even though the exact worst-case query depth remains all 16 bits.

\paragraph{Full $D(4)$ as a cube edge-cover function.}
Each genuine 3-clause on four variables falsifies exactly two assignments, differing in the omitted variable. Hence the 32 clause bits are naturally the 32 edges of $Q_4$. A selected formula is UNSAT exactly when every assignment is falsified by at least one selected clause, equivalently when the selected edges form an edge cover of all 16 vertices of $Q_4$.

The $Q_4$ automorphism group acts transitively on its 32 edges: translations move base vertices and coordinate permutations move edge directions. Edge-cover is invariant under this action. Since
\[
32=2^5
\]
and the all-zero and all-one edge sets have different outputs, Rivest--Vuillemin's prime-power transitivity theorem gives
\[
\boxed{D(f)=32}.
\]
Therefore the full four-variable formula-presence SAT function is evasive: every deterministic adaptive query algorithm has some input on which it must inspect all 32 distinct clause bits.

The 2432-node FBDD was obtained by allowing state-dependent deviations among the first two relevant variables of the best sign-major priority. One million seeded random presence vectors were checked against direct SAT evaluation with zero mismatches. This is validation of the implementation, not proof of size optimality.

\paragraph{Time versus state.}
The chain control shows
\[
\text{DAG size can collapse while worst-case query depth remains unchanged}.
\]
For $D(4)$, evasiveness fixes the worst-case distinct-query depth at 32 even if rereads are later allowed. Thus a read-twice or RECOMP improvement can only come from a STORE/FORGET/RECOMPUTE state-size tradeoff; it cannot reduce the number of distinct clause bits that some computation must inspect.

\paragraph{Literature boundary.}
FBDDs are the adaptive-order read-once branching-program model. Exact minimization is nontrivial, and lower-bound methods for read-once branching programs use subfunction, communication, rectangle, rank, and related combinatorial measures. The literature on read-$k$ and $(1,+k)$ branching programs shows that even a small allowance for rereads moves to a substantially harder lower-bound regime.

\paragraph{Barrier.}
Evasiveness is only a linear input-query lower bound. Neither the depth-32 theorem nor the finite FBDD sizes imply superpolynomial SAT time. No theorem here forces an arbitrary polynomial-time RAM/Turing SAT algorithm to instantiate the same branching program or state resource.

\paragraph{Next target.}
Use the $Q_4$ edge-cover formulation to synthesize read-twice or $(1,+k)$ branching programs and search for a size reduction below 2432. Keep the exact depth-32 theorem fixed while measuring how much state can be traded for legal recomputation.

\begin{thebibliography}{9}
\bibitem{RV76}
R.~L.~Rivest and J.~Vuillemin,
\emph{On Recognizing Graph Properties from Adjacency Matrices},
Theoretical Computer Science 3(3):371--384, 1976.
doi:10.1016/0304-3975(76)90053-0.

\bibitem{LiMcKenzie24}
Y.~Li and P.~McKenzie,
\emph{Perspective on Complexity Measures Targeting Read-Once Branching Programs},
Information and Computation 301(A):105230, 2024.
doi:10.1016/j.ic.2024.105230.

\bibitem{GuentherDrechsler02}
W.~G\"unther and R.~Drechsler,
\emph{Minimization of Free BDDs},
Integration 32(1--2):41--59, 2002.
doi:10.1016/S0167-9260(02)00041-X.

\bibitem{SavickyZak00}
P.~Savick\'y and S.~\v{Z}\'ak,
\emph{A Read-Once Lower Bound and a $(1,+k)$-Hierarchy for Branching Programs},
Theoretical Computer Science 238(1--2):347--362, 2000.
doi:10.1016/S0304-3975(98)00219-9.
\end{thebibliography}

% END_OF_NOTE::ADAPTIVE_FORMULA_FBDD5
